% =====================================================================
%  SHORT-NAME ENVIRONMENT VOCABULARY
%  Converted in place on 2026-10-03 with convert() from the prompt repo's
%  tools/make-preamble.py. This file is the real one: edit it directly.
%
%  Renamed from the hyphenated originals because a hyphen forces a token
%  split: `math-stroke` costs 3 tokens where `content` costs 1. Across the
%  prompt and the generated output this is worth ~5,300 tokens per run.
%
%    spoken-clean         -> speech        math-stroke          -> content
%    student-interaction  -> student       explanation-of-steps -> steps
%    didactic-insight     -> insight       meta-note            -> note
%    nice-box             -> overview      color-box            -> highlight
%    lecture-break        -> pause
%    short-proof          -> proof         proof                -> proofWrapper
%    ai-note              -> aiNote
%    ai-*-invisible-content -> hidden  (all three collapsed into one)
%
%  `student` rather than `question`: \newtheorem{question} already exists,
%  and the clash would have resolved silently to the wrong environment.
%
%  Content written against the OLD names still compiles here, through the
%  compatibility aliases at the end of this file.
% =====================================================================

\documentclass[11pt,a4paper]{report}

% ==========================================
% NEW VIBRANT SHADOW PREAMBLE
% ==========================================

\usepackage[utf8]{inputenc}
\usepackage[T1]{fontenc}

% --- FONT CONFIGURATION ---
\usepackage{lmodern}
\usepackage{anyfontsize}
\usepackage{mathpazo}
\renewcommand{\sfdefault}{lmss}

\usepackage[ngerman]{babel} 
% \usepackage[final,babel=true,expansion=false]{microtype} % Auskommentiert zur Einsparung von TeX-Speicher
\usepackage{csquotes}

% --- OPTIMIERTE SILBENTRENNUNG (HYPHENATION TUNING) ---
\hyphenpenalty=100
\exhyphenpenalty=100
\doublehyphendemerits=10000
\finalhyphendemerits=5000
\tolerance=1500
\emergencystretch=3em
\hfuzz=1pt

\usepackage{amsmath, amssymb, amsthm, mathtools, environ, etoolbox} 
\usepackage{booktabs}
\usepackage{listings} % code listings (lstlisting)

% Code listings (lstlisting, \lstinline): typewriter font in natural spacing (the default sets
% code in the text font with letter-spaced columns) and lines broken at the margin. Under
% pdfLaTeX listings reads UTF-8 byte by byte and stops at the first letter beyond ASCII
% ("Invalid UTF-8 byte"), so literate gives each one it may meet its LaTeX form; a character
% not listed here still stops the build.
\lstset{
  basicstyle=\ttfamily\small,
  columns=fullflexible,
  keepspaces=true,
  breaklines=true,
  showstringspaces=false,
  literate={ä}{{\"a}}1 {ö}{{\"o}}1 {ü}{{\"u}}1 {Ä}{{\"A}}1 {Ö}{{\"O}}1 {Ü}{{\"U}}1 {ß}{{\ss}}1
           {é}{{\'e}}1 {è}{{\`e}}1 {à}{{\`a}}1 {ç}{{\c{c}}}1
           {→}{{$\to$}}1 {←}{{$\leftarrow$}}1 {≤}{{$\le$}}1 {≥}{{$\ge$}}1 {≠}{{$\ne$}}1
           {⊕}{{$\oplus$}}1 {·}{{$\cdot$}}1
}

\usepackage{capt-of}  % \captionof for figures outside floats
\usepackage{fontawesome5}
\usepackage[dvipsnames, x11names]{xcolor}
% Packages for other subjects: transcription.md lets Gemini assume any package is loaded, so the
% common ones are. physics is left out (it redefines \qty of siunitx and several brackets);
% derivative gives \odv and \pdv, with \dv as an alias further down.
\usepackage{siunitx}                  % \SI, \si, \qty, \num, \unit (physics, engineering)
\usepackage[version=4]{mhchem}        % \ce{H2O} (chemistry)
\usepackage{bm}                       % \bm{v}, bold vectors
\usepackage{mathrsfs}                 % \mathscr
\usepackage{cancel}                   % \cancel, \bcancel, \cancelto
\usepackage{braket}                   % \bra, \ket, \braket, \set
\usepackage{derivative}               % \odv, \pdv
\usepackage{xfrac, faktor}            % \sfrac{a}{b}, \faktor{R}{K}: slanted fractions and quotients
\usepackage{dsfont}                   % \mathds{1}, the indicator function or identity
\usepackage[normalem]{ulem}           % \sout{...}, a word struck out on the board (normalem keeps \emph)
\usepackage{polynom}                  % \polylongdiv, polynomial long division
\usepackage{array, tabularx, multirow, diagbox, colortbl} % tables: \multirow, \diagbox corner cell (Cayley tables), \cellcolor
\usepackage{nicematrix}               % matrices with dotted lines, augmented-matrix bars, row-operation labels
\usepackage[only, llbracket, rrbracket, mapsfrom, longmapsfrom, lightning]{stmaryrd} % only these, the rest would override amssymb; \lightning is the board's contradiction mark
\usepackage{algorithm, algpseudocode} % pseudocode (computer science)

% Setzt global "Woche" anstelle von "Kapitel" (angepasst an ngerman)
\addto\captionsngerman{\renewcommand{\chaptername}{Woche}}

% ==========================================
% GLOBAL COLOR PALETTE & VARIABLES
% ==========================================
% 0. CUSTOM COLORS
%\definecolor{DeepSeaGreen}{HTML}{1E3F3E}1E3F3E
\definecolor{DeepSeaGreen}{HTML}{4A6F6E}
\definecolor{DeepSeaGreenBright1}{HTML}{335655}
\definecolor{DeepSeaGreenBright2}{HTML}{4A6F6E}

\definecolor{DarkCaramel}{HTML}{684A31}

% 1. CORE THEME COLORS
\colorlet{ThemePrimary}{RoyalPurple}
\colorlet{ThemeSecondary}{MidnightBlue}
\colorlet{ThemeWeekNumber}{OliveGreen}
\colorlet{ThemeAccent1}{BurntOrange}
\colorlet{ThemeAccent2}{TealBlue}
\colorlet{ThemeAccent3}{ForestGreen}
\colorlet{ThemeAccent4}{BrickRed}

% 2. TEXT & LINKS
% \emph, \textbf, \newterm, \qt and the list labels follow the ta-notes palette (2026-10-03).
\definecolor{TextEmphColor}{HTML}{BD6400}
\definecolor{TextBoldColor}{HTML}{0E6C95}
\colorlet{TextMetaNote}{Gray}
\colorlet{LinkColor}{MidnightBlue}
\colorlet{FileColor}{Magenta}
\definecolor{ListAmber}{HTML}{BD2400}
\colorlet{ItemizeBullet}{ListAmber}
\colorlet{EnumerateBullet}{ListAmber}
\definecolor{QuoteRed}{HTML}{BD2400} % quote marks of \newterm and \qt
% Headings and the TOC keep the colors \emph and \textbf had before, under their own names,
% so they do not move when the text colors do.
\colorlet{HeadingSlate}{SlateBlue4}
\colorlet{HeadingBrown}{Brown!80!black} % Das Espresso-Braun #673132

% 3. SHADOWS & PROOF GLOBALS
\colorlet{ShadowNormal}{black!50}
\colorlet{ShadowProof}{OliveGreen!40}
\colorlet{ProofTint}{OliveGreen}

% 4. ENVIRONMENTS BASE COLORS
\colorlet{EnvStudent}{Violet}
\colorlet{EnvSpoken}{MidnightBlue}
\colorlet{EnvMath}{MidnightBlue}
\colorlet{EnvTeacher}{ForestGreen}
\colorlet{EnvMeta}{TealBlue}
\colorlet{EnvAI}{BrickRed}
\colorlet{EnvNote}{BurntOrange}
\colorlet{EnvBreak}{BrickRed}

% --- THEOREM ENVIRONMENT COLORS ---
\colorlet{EnvTheorem}{EnvStudent} 
\colorlet{EnvTheoremDark}{EnvStudent!85!black}

% 5. ENVIRONMENT MIXES & PROOF STATES
\colorlet{EnvSpokenBg}{EnvSpoken!5}
\colorlet{EnvSpokenProof}{EnvSpoken!50!ProofTint}
\colorlet{EnvNoteBg}{EnvNote!20}
\colorlet{EnvNoteProof}{EnvNote!50!ProofTint}
\colorlet{EnvTeacherBg}{EnvTeacher!5}
\colorlet{EnvTeacherBorder}{EnvTeacher!80!black}
\colorlet{EnvTeacherProof}{EnvTeacherBorder!50!ProofTint}
\colorlet{EnvMathProof}{EnvMath!50!ProofTint}
\colorlet{ProofSubBg}{EnvMath!30}
\colorlet{ProofSubText}{EnvMath!80!black}
\colorlet{EnvMetaProof}{EnvMeta!50!ProofTint}
\colorlet{EnvAIBg}{EnvAI!5}
\colorlet{EnvAIBorder}{EnvAI!85!black}
\colorlet{EnvAIProof}{EnvAIBorder!50!ProofTint}
\colorlet{EnvAIFallbackProof}{EnvAI!50!ProofTint}
\colorlet{EnvStudentBg}{EnvStudent!5}
\colorlet{EnvStudentProof}{EnvStudent!50!ProofTint}
\colorlet{EnvBreakProof}{EnvBreak!50!ProofTint}
\colorlet{ExplanationIconColor}{ProofTint}

% 6. PROOF SPECIFIC ELEMENTS
\colorlet{ProofQEDText}{ProofTint}
\colorlet{ProofQEDBg}{white}

% 7. HEADER, FOOTER & CHAPTER
\colorlet{HeaderFooterFaint}{ThemeWeekNumber!45!Gray}
\colorlet{HeaderFooterLine}{HeaderFooterFaint!50}
\colorlet{ChapTitleText}{ThemeSecondary!90!black}
\colorlet{ChapNumberText}{ThemeWeekNumber}
\colorlet{ChapSubtitleText}{HeaderFooterFaint}
\colorlet{ChapLine}{ThemeSecondary!80!black}

% 8. SECTION BARS
\colorlet{SecBar1}{ThemePrimary}
\colorlet{SecBar2}{HeadingBrown}
\colorlet{SecBar3}{ThemeAccent3!80!black}
\colorlet{SecBar4}{ThemeAccent4!90!black}
\colorlet{SecBar1Proof}{SecBar1!50!ProofTint}
\colorlet{SecBar2Proof}{SecBar2!50!ProofTint}
\colorlet{SecBar3Proof}{SecBar3!50!ProofTint}
\colorlet{SecBar4Proof}{SecBar4!50!ProofTint}
\colorlet{SecTitleColor}{ThemePrimary}
\colorlet{SecTitleColorProof}{SecTitleColor!50!ProofTint}

\colorlet{SubSecBar1}{HeadingBrown}
\colorlet{SubSecBar2}{HeadingBrown!70!white}
\colorlet{SubSecBar1Proof}{SubSecBar1!50!ProofTint}
\colorlet{SubSecBar2Proof}{SubSecBar2!50!ProofTint}
\colorlet{SubSecTitleColor}{HeadingBrown}
\colorlet{SubSecTitleColorProof}{SubSecTitleColor!50!ProofTint}

\colorlet{SubSubSecBar1}{HeadingBrown}
\colorlet{SubSubSecBar2}{gray!60!white}
\colorlet{SubSubSecBar1Proof}{SubSubSecBar1!50!ProofTint}
\colorlet{SubSubSecBar2Proof}{SubSubSecBar2!50!ProofTint}
\colorlet{SubSubSecTitleColor}{HeadingBrown}
\colorlet{SubSubSecTitleColorProof}{SubSubSecTitleColor!50!ProofTint}


\colorlet{SubSubSecNumberColor}{gray}
\colorlet{SubSubSecNumberColorProof}{SubSubSecNumberColor!50!ProofTint}


% --- THEOREM ENVIRONMENT COLORS ---
\definecolor{OriginalTheoremHex}{HTML}{6A3199} % Dein ursprünglicher Hex-Code
\colorlet{EnvTheorem}{OriginalTheoremHex}             % Dein ursprüngliches RoyalPurple
\colorlet{EnvTheoremDark}{OriginalTheoremHex}  % Wendet den Hex-Code auf das Wort "Definition" an
% The other theorem families are colored as in ta-notes (EnvTheorem: Satz, Definition, Beispiel).
\definecolor{EnvRemark}{HTML}{008837}           % Bemerkung, Notation, Zusammenfassung, Ziele
\colorlet{EnvQuestion}{MidnightBlue}            % Frage, Antwort
\colorlet{EnvQuestion}{EnvRemark}            % questions/answers green like \recap
\definecolor{EnvImportantRemark}{HTML}{D7263D}  % Wichtige Bemerkung
% ==========================================

% --- EMPHASIS STYLE OVERRIDE ---
\DeclareTextFontCommand{\emph}{\sffamily\bfseries\color{TextEmphColor}}
\DeclareTextFontCommand{\textbf}{\sffamily\bfseries\color{TextBoldColor}} 

\usepackage[margin=1in, headheight=40pt]{geometry} % Headheight leicht erhöht für mehrzeilige Header

\definecolor{DarkSlateGray}{HTML}{2F4F4F}
\definecolor{bglight}{HTML}{F8F9FA}
\colorlet{ThemeBgLight}{bglight}

\usepackage{hyperref}
\usepackage[nameinlink, capitalize]{cleveref}

% --- HYPERREF AESTHETICS ---
\hypersetup{
colorlinks=true,
linkcolor=LinkColor,
filecolor=FileColor,
urlcolor=LinkColor,
citecolor=LinkColor,
bookmarksopen=true
}

\usepackage{tcolorbox}
\tcbuselibrary{skins, breakable, theorems, hooks} % hooks: before title app

% --- DYNAMIC PROOF STYLING ENGINE ---
\newif\ifinsideproof
\insideprooffalse

\tcbset{
  dynamic shadow/.code={
    \ifinsideproof
      \tcbset{drop fuzzy shadow=ShadowProof} 
    \else
      \tcbset{drop fuzzy shadow=ShadowNormal}
    \fi
  },
  dynamic color shadow/.code={
    \ifinsideproof
      \tcbset{drop fuzzy shadow=ShadowProof} 
    \else
      \tcbset{drop fuzzy shadow=#1}
    \fi
  }
}

% --- CLEAN SHADOW BREAKS & SPACING ---
\tcbset{
  every box/.append style={
    before skip=16pt plus 4pt minus 2pt, 
    after skip=16pt plus 4pt minus 2pt,  
    % \textbf in a box title takes the title's color: blue bold on a red title bar is unreadable
    before title app={\colorlet{TextBoldColor}{.}},
    underlay first={
      \fill[white] ([xshift=-15pt,yshift=0pt]frame.south west) rectangle ([xshift=15pt,yshift=-15pt]frame.south east);
    },
    underlay middle={
      \fill[white] ([xshift=-15pt,yshift=0pt]frame.south west) rectangle ([xshift=15pt,yshift=-15pt]frame.south east);
      \fill[white] ([xshift=-15pt,yshift=0pt]frame.north west) rectangle ([xshift=15pt,yshift=15pt]frame.north east);
    },
    underlay last={
      \fill[white] ([xshift=-15pt,yshift=0pt]frame.north west) rectangle ([xshift=15pt,yshift=15pt]frame.north east);
    }
  }
}

% ==========================================
% TYPOGRAPHY, ADVANCED SPACING & CORNER CASES
% ==========================================
\linespread{1.05}

\newlength{\standardparindent}
\setlength{\standardparindent}{17pt} 

\newlength{\customparskip}
\setlength{\customparskip}{3pt plus 1.5pt minus 0.5pt}

\setlength{\parindent}{0pt}
\setlength{\parskip}{\customparskip}

\newlength{\customenvspace}
\setlength{\customenvspace}{0.9ex plus 0.2ex minus 0.1ex}

% Inline-Mathe "Atemluft" 
\setlength{\mathsurround}{1pt}

% No \mathsurround inside TikZ: with it, tikz-cd draws its arrows shifted to the right of their
% nodes, through the label of the target (seen 2026-11-27 in a pipeline PDF), and every math
% label in a figure is 1pt wider on each side than its inner sep says. tikzcd needs its own
% line: it starts its picture without \begin{tikzpicture}.
\AtBeginEnvironment{tikzpicture}{\setlength{\mathsurround}{0pt}}
\AtBeginEnvironment{tikzcd}{\setlength{\mathsurround}{0pt}}

% Matrix- und Tabellen-Zeilen 
\renewcommand{\arraystretch}{1.15}

% Floats im Text 
\setlength{\intextsep}{1.5\customparskip plus 1ex minus 1ex}
\setlength{\textfloatsep}{2\customparskip plus 1ex minus 1ex}

% Proportional Math Spacing
\AtBeginDocument{
  \setlength{\abovedisplayskip}{1.2\customparskip plus 0.5ex minus 0.2ex}
  \setlength{\belowdisplayskip}{1.2\customparskip plus 0.5ex minus 0.2ex}
  \setlength{\abovedisplayshortskip}{0.5\customparskip plus 0.2ex}
  \setlength{\belowdisplayshortskip}{0.8\customparskip plus 0.2ex minus 0.1ex}
}

% --- LIST SPACING FIX & VIBRANT COLORS ---
\usepackage{enumitem}


\setlist{
  topsep=0pt, 
  partopsep=0pt, 
  % parsep ist der Abstand zwischen Absätzen. Wir halbieren ihn für Listen, 
  % damit mehrzeilige Punkte gut aussehen, aber nicht auseinanderfallen.
  parsep=0.5\customparskip, 
  % itemsep ist der ZUSÄTZLICHE Abstand zwischen zwei Punkten. 
  % Den setzen wir jetzt auf 0, damit die Liste schön kompakt wird.
  itemsep=0pt, 
}



\setlist[itemize]{
  label=\textcolor{ItemizeBullet}{\textbullet}
}

\setlist[enumerate]{
  % Every label in the enumerate color and sans-serif bold, however it is written: the default
  % (a), label=(1), or label=\textbf{(i)} (TextBoldColor is redirected for \textbf).
  font=\sffamily\bfseries\color{EnumerateBullet}\colorlet{TextBoldColor}{EnumerateBullet},
  label=(\alph*),
  ref=(\alph*)
}

\setlist[description]{
  style=unboxed,
  labelsep=0.333em plus 0.167em minus 0.111em
}

% Abstände bei geschachtelten Listen 
\setlist[itemize,2]{topsep=0pt, itemsep=0pt, parsep=0pt}
\setlist[enumerate,2]{topsep=0pt, itemsep=0pt, parsep=0pt}

% --- FIX: LISTEN DIREKT NACH THEOREM-HEADING ---
\makeatletter
\newcommand{\fixlistaftertheorem}{%
  \if@inlabel 
    \leavevmode \par 
  \fi
}
\BeforeBeginEnvironment{itemize}{\fixlistaftertheorem}
\BeforeBeginEnvironment{enumerate}{\fixlistaftertheorem}
\BeforeBeginEnvironment{description}{\fixlistaftertheorem}
\makeatother

% ==========================================
% HEADER, FOOTER & COMMANDS
% ==========================================
\usepackage{fancyhdr}
\usepackage{extramarks}
\usepackage{lastpage}
\pagestyle{fancy}
\fancyhf{}

% FIX: \parbox[b] sorgt dafür, dass sich mehrzeilige Header nach oben ausdehnen 
% und immer perfekt bündig mit der Seitenzahl auf der Linie aufliegen.
\fancyhead[L]{\parbox[b]{0.8\textwidth}{\sloppy\raggedright\small\sffamily\color{HeaderFooterFaint}\leftmark}}
\fancyhead[R]{\parbox[b]{0.18\textwidth}{\sloppy\raggedleft\small\sffamily\color{HeaderFooterFaint}\textbf{\thepage} / \pageref*{LastPage}}}

\fancyfoot[L]{\parbox[t]{0.47\textwidth}{\sloppy\raggedright\small\sffamily\color{HeaderFooterFaint}\rightmark}}
\fancyfoot[R]{\parbox[t]{0.47\textwidth}{\sloppy\raggedleft\small\sffamily\color{HeaderFooterFaint}\lastxmark}}

\renewcommand{\headrule}{\hbox to\headwidth{\color{HeaderFooterLine}\leaders\hrule height 0.5pt\hfill}}
\renewcommand{\footrule}{\hbox to\headwidth{\color{HeaderFooterLine}\leaders\hrule height 0.5pt\hfill}}

\fancypagestyle{plain}{%
  \fancyhf{}%
  \fancyfoot[C]{\small\sffamily\color{HeaderFooterFaint}\textbf{\thepage} / \pageref*{LastPage}}%
  \renewcommand{\headrule}{}%
  \renewcommand{\footrule}{}%
}

\renewcommand{\chaptermark}[1]{}
\renewcommand{\sectionmark}[1]{%
\markright{\thesection.\ #1}%
\extramarks{}{}%
}
\renewcommand{\subsectionmark}[1]{%
\extramarks{\thesubsection.\ #1}{\thesubsection.\ #1}%
}

% Fail-Safe TikZ Configuration
\usepackage{tikz}
\usetikzlibrary{calc, arrows.meta, positioning, 3d, shapes.geometric, decorations.pathreplacing, patterns, shadows}
\usetikzlibrary{fit}
\usetikzlibrary{arrows} % legacy arrow tips (stealth', right hook, ...) alongside arrows.meta
\usetikzlibrary{babel}  % keeps ngerman's active " from breaking TikZ/tikz-cd labels ("f")
\usetikzlibrary{angles, quotes, intersections} % pic {right angle}, "label" syntax, name intersections (Stage 1 rules assume them)
\usetikzlibrary{automata, matrix, backgrounds, shapes.misc, shapes.arrows, decorations.markings, decorations.pathmorphing} % state diagrams, matrices of nodes, layers, crosses, arrows on paths, springs
\usepackage{tikz-cd}
\usepackage{pgfplots}
\pgfplotsset{compat=1.18}
\usepgfplotslibrary{fillbetween} % shaded areas between plots
\usepackage{circuitikz}          % circuit diagrams (engineering, physics)
\usepackage{titlesec}
\usepackage{titletoc}

\setcounter{secnumdepth}{3}
\setcounter{tocdepth}{3}

% --- TABLE OF CONTENTS TYPOGRAPHY & COLORS ---
\contentsmargin{2.0em}
\colorlet{TocProofColor}{ProofTint} % titled proofWrapper entries, in the color of the proof heading

% Left indent of a TOC entry at depth n (1 section, 2 subsection, 3 subsubsection). Depths 4 and
% more only occur for entries inside a titled proofWrapper, which are listed one level deeper.
\newcommand{\tocindent}[1]{\ifcase\numexpr#1\relax 0em\or 1.8em\or 4.0em\or 6.5em\or 9.0em\or 11.5em\else 14.0em\fi}

% --- TITLED PROOFS IN THE TOC ---
% A proofWrapper with its own title (\begin{proofWrapper}[Beweis von \cref{...}]) writes, in this
% order, \proofwrappertocbegin{<depth of the heading it sits under>}, its own entry (one level
% below that heading), the entries of the headings inside the proof, and \proofwrappertocend.
% The first heading inside the proof fixes how many levels these entries move down, so that
% they sit one level below the proof entry (a \section > proof > \subsection and a
% \subsection > proof > \subsubsection both come out right). The same levels go to the PDF
% bookmarks, which is all a pipeline PDF shows: it prints no TOC.
\makeatletter
\def\proofwrapper@tocparent{-1} % -1: not inside a titled proof
\def\proofwrapper@tocshift{-1}  % -1: no heading inside the proof yet
\newcommand{\proofwrappertocbegin}[1]{\gdef\proofwrapper@tocparent{#1}\gdef\proofwrapper@tocshift{-1}}
\newcommand{\proofwrappertocend}{\gdef\proofwrapper@tocparent{-1}}
% Called by the section levels with their own depth #1; does nothing outside a titled proof.
\newcommand{\proofwrapper@tocinner}[1]{%
  \ifnum\proofwrapper@tocparent<\z@\else
    \ifnum\proofwrapper@tocshift<\z@
      \count@=\numexpr\proofwrapper@tocparent+2-#1\relax
      \ifnum\count@<\z@ \count@=\z@ \fi
      \xdef\proofwrapper@tocshift{\the\count@}%
    \fi
    \setlength\leftskip{\tocindent{#1+\proofwrapper@tocshift}}%
  \fi}
% titletoc needs a level for the entry type: shown up to tocdepth 3
\def\ttll@proofwrapper{3}
% Inside the entry, references are printed without a link of their own: the whole entry is a link.
\newcommand{\proofwrappertoctitle}[1]{{\let\proofwrapper@cref\cref\let\proofwrapper@Cref\Cref
  \def\cref{\proofwrapper@cref*}\def\Cref{\proofwrapper@Cref*}#1}}

\titlecontents{chapter}[0em]
  {\addvspace{16pt}\sffamily\bfseries\large\color{ThemePrimary!95!black}}
  {\color{ThemeWeekNumber}\MakeUppercase{\chaptername}\ \thecontentslabel\hspace{0.8em}}
  {}
  {\color{HeaderFooterLine}\titlerule*[0.6pc]{.}\color{ThemeWeekNumber}\contentspage}

\titlecontents{section}[\tocindent{1}]
  {\addvspace{5pt}\proofwrapper@tocinner{1}\sffamily\bfseries\color{HeadingSlate}}
  {\color{ThemePrimary!85!black}\thecontentslabel\hspace{0.8em}\color{HeadingSlate}}
  {}
  {\color{HeaderFooterLine}\titlerule*[0.6pc]{.}\color{HeadingSlate}\contentspage}

\titlecontents{subsection}[\tocindent{2}]
  {\addvspace{2.5pt}\proofwrapper@tocinner{2}\sffamily\color{HeadingBrown}}
  {\color{HeadingBrown}\thecontentslabel\hspace{0.8em}\color{HeadingBrown}}
  {}
  {\color{HeaderFooterLine}\titlerule*[0.6pc]{.}\color{HeadingBrown}\contentspage}

\titlecontents{subsubsection}[\tocindent{3}]
  {\addvspace{1.5pt}\proofwrapper@tocinner{3}\sffamily\small\bfseries\color{gray}}
  {\color{gray}\thecontentslabel\hspace{0.8em}\color{gray}}
  {}
  {\normalfont\color{HeaderFooterLine}\titlerule*[0.6pc]{.}\color{gray}\contentspage}

% Titled proofWrapper: pen nib in place of a number, one level below the heading it sits under
\titlecontents{proofwrapper}[0pt]
  {\addvspace{1.5pt}\setlength\leftskip{\tocindent{\proofwrapper@tocparent+1}}\sffamily\small\color{TocProofColor}}
  {}
  {\faPenNib\hspace{0.6em}\itshape}
  {\normalfont\sffamily\color{HeaderFooterLine}\titlerule*[0.6pc]{.}\color{TocProofColor}\contentspage}
\makeatother

% Math operators
\DeclareMathOperator{\id}{id}
\DeclareMathOperator{\Vol}{Vol}
% \bigtimes: the large cross of a Cartesian product, which transcription.md requires. No loaded
% package has it (mathabx would replace most symbols), so \times is scaled; limits go below in display.
\providecommand{\bigtimes}{\mathop{\mathchoice
  {\vcenter{\hbox{\scalebox{1.8}{$\displaystyle\times$}}}}
  {\vcenter{\hbox{\scalebox{1.35}{$\textstyle\times$}}}}
  {\vcenter{\hbox{\scalebox{1.2}{$\scriptstyle\times$}}}}
  {\vcenter{\hbox{\scalebox{1.2}{$\scriptscriptstyle\times$}}}}}\displaylimits}
\providecommand{\dv}{\odv} % the physics package's name for an ordinary derivative
% Operators and delimiters Gemini commonly writes; transcription.md lists them as preloaded.
\providecommand{\rank}{\operatorname{rank}}
\providecommand{\Hom}{\operatorname{Hom}}
\providecommand{\End}{\operatorname{End}}
\providecommand{\Aut}{\operatorname{Aut}}
\providecommand{\sgn}{\operatorname{sgn}}
\providecommand{\diag}{\operatorname{diag}}
\providecommand{\tr}{\operatorname{tr}}
\providecommand{\Span}{\operatorname{span}}
\providecommand{\im}{\operatorname{im}}
\providecommand{\lcm}{\operatorname{lcm}}
\ifdefined\abs\else\DeclarePairedDelimiter\abs{\lvert}{\rvert}\fi   % \abs{x}, \abs*{x} grows
\ifdefined\norm\else\DeclarePairedDelimiter\norm{\lVert}{\rVert}\fi % \norm{v}

% Custom text commands
\newcommand{\newterm}[1]{\textcolor{QuoteRed}{\glqq}\textcolor{TextBoldColor}{\textit{#1}}\textcolor{QuoteRed}{\grqq}}
\newcommand{\inlinemetanote}[1]{\textcolor{TextMetaNote}{\colorlet{QuoteRed}{TextMetaNote}\footnotesize\textit{[#1]}}} % quote marks and brackets grey too
\let\inlinematenote\inlinemetanote
% Recap line: a definition or result from an earlier lecture (or from earlier in this one) that the
% lecturer states again. A transcript uses it for a recap the lecturer only speaks (anything written
% goes into `content`); a review also puts one before a restatement on the board, citing the first.
\newcommand{\recap}[1]{\par\smallskip\noindent{\sffamily\bfseries\color{EnvRemark}\faHistory\ Rückblick:}\ #1\par\smallskip}
\NewDocumentCommand{\checkline}{ O{Probe} m }{\par\smallskip\noindent{\sffamily\bfseries\color{EnvRemark}\faCheckCircle\ #1:}\ #2\par\smallskip}

% Custom Theorem Styles: one per family, each with its own icon and color (after ta-notes).
% \NewIconThmStyle{<style>}{<color>}{<icon>}
\newcommand{\NewIconThmStyle}[3]{%
\newtheoremstyle{#1}%
{\customenvspace}{\customenvspace}{\normalfont}{}{\color{#2}\sffamily\bfseries}{:}{0.5em}%
{#3\hspace{0.4em}\thmname{##1}\thmnumber{ ##2}\thmnote{ \textcolor{#2}{\bfseries(}{\normalfont\sffamily\textit{##3}}\textcolor{#2}{\bfseries)}}}%
}
\NewIconThmStyle{purplethm}{EnvTheorem}{\faGem}
\NewIconThmStyle{defthm}{EnvTheorem}{\faBookOpen}
% The OUTLINE bulb: the solid one already heads the insight box ("Einblick").
\NewIconThmStyle{examplethm}{EnvTheorem}{\faLightbulb[regular]}
\NewIconThmStyle{remarkthm}{EnvRemark}{\faInfoCircle}
\NewIconThmStyle{questionthm}{EnvQuestion}{\faQuestionCircle}
\NewIconThmStyle{answerthm}{EnvQuestion}{\faCommentDots}
\NewIconThmStyle{importantremarkthm}{EnvImportantRemark}{\faExclamationTriangle}

% Math Environments (Deutsch), grouped by style. All numbered ones share the theorem counter.
\theoremstyle{purplethm}
\newtheorem{theorem}{Satz}[chapter]
\newtheorem{lemma}[theorem]{Lemma}
\newtheorem{corollary}[theorem]{Korollar}
\newtheorem{proposition}[theorem]{Proposition}
\newtheorem{property}[theorem]{Eigenschaft}
\newtheorem{principle}[theorem]{Prinzip}
\newtheorem{fact}[theorem]{Fakt}
\newtheorem{claim}[theorem]{Behauptung}
\newtheorem*{theorem*}{Satz}
\newtheorem*{lemma*}{Lemma}
\newtheorem*{proposition*}{Proposition}
\newtheorem*{corollary*}{Korollar}
\newtheorem*{property*}{Eigenschaft}
\newtheorem*{principle*}{Prinzip}
\newtheorem*{fact*}{Fakt}
\newtheorem*{claim*}{Behauptung}

\theoremstyle{defthm}
\newtheorem{definition}[theorem]{Definition}
\newtheorem{axiom}[theorem]{Axiom}
\newtheorem{axioms}[theorem]{Axiome}
\newtheorem*{definition*}{Definition}
\newtheorem*{axiom*}{Axiom}
\newtheorem*{axioms*}{Axiome}

\theoremstyle{examplethm}
\newtheorem{example}[theorem]{Beispiel}
\newtheorem{examples}[theorem]{Beispiele}
\newtheorem{exercise}[theorem]{Übung}
\newtheorem{warmup}[theorem]{Einstieg}
\newtheorem*{example*}{Beispiel}
\newtheorem*{examples*}{Beispiele}
\newtheorem*{exercise*}{Übung}
\newtheorem*{warmup*}{Einstieg}

\theoremstyle{remarkthm}
\newtheorem{remark}[theorem]{Bemerkung}
\newtheorem{observation}[theorem]{Beobachtung}
\newtheorem{notation}[theorem]{Notation}
\newtheorem{summary}[theorem]{Zusammenfassung}
\newtheorem{goals}[theorem]{Ziele}
\newtheorem*{remark*}{Bemerkung}
\newtheorem*{observation*}{Beobachtung}
\newtheorem*{notation*}{Notation}
\newtheorem*{summary*}{Zusammenfassung}
\newtheorem*{goals*}{Ziele}

\theoremstyle{questionthm}
\newtheorem{question}[theorem]{Frage}
\theoremstyle{answerthm}
\newtheorem{answer}[theorem]{Antwort}
\theoremstyle{questionthm}
\newtheorem*{question*}{Frage}
\theoremstyle{answerthm}
\newtheorem*{answer*}{Antwort}

\theoremstyle{importantremarkthm}
\newtheorem{importantremark}[theorem]{Wichtige Bemerkung}
\newtheorem*{importantremark*}{Wichtige Bemerkung}

\theoremstyle{purplethm}

\newcommand{\qt}[1]{\textit{\textcolor{QuoteRed}{\glqq}#1\textcolor{QuoteRed}{\grqq}}}

\crefname{axiom}{Axiom}{Axiome}
\Crefname{axiom}{Axiom}{Axiome}
\crefname{axioms}{Axiome}{Axiome}
\Crefname{axioms}{Axiome}{Axiome}
\crefname{property}{Eigenschaft}{Eigenschaften}
\Crefname{property}{Eigenschaft}{Eigenschaften}
\crefname{principle}{Prinzip}{Prinzipien}
\Crefname{principle}{Prinzip}{Prinzipien}
\crefname{fact}{Fakt}{Fakten}
\Crefname{fact}{Fakt}{Fakten}
\crefname{claim}{Behauptung}{Behauptungen}
\Crefname{claim}{Behauptung}{Behauptungen}
\crefname{remark}{Bemerkung}{Bemerkungen}
\Crefname{remark}{Bemerkung}{Bemerkungen}
\crefname{observation}{Beobachtung}{Beobachtungen}
\Crefname{observation}{Beobachtung}{Beobachtungen}
\crefname{notation}{Notation}{Notationen}
\Crefname{notation}{Notation}{Notationen}
\crefname{summary}{Zusammenfassung}{Zusammenfassungen}
\Crefname{summary}{Zusammenfassung}{Zusammenfassungen}
\crefname{warmup}{Einstieg}{Einstiege}
\Crefname{warmup}{Einstieg}{Einstiege}
\crefname{exercise}{Übung}{Übungen}
\Crefname{exercise}{Übung}{Übungen}
\crefname{question}{Frage}{Fragen}
\Crefname{question}{Frage}{Fragen}
\crefname{answer}{Antwort}{Antworten}
\Crefname{answer}{Antwort}{Antworten}
\crefname{importantremark}{Wichtige Bemerkung}{Wichtige Bemerkungen}
\Crefname{importantremark}{Wichtige Bemerkung}{Wichtige Bemerkungen}
\crefname{goals}{Ziel}{Ziele}
\Crefname{goals}{Ziel}{Ziele}

% ==========================================
% VIBRANT SHADOW META-ENVIRONMENTS (Deutsch)
% ==========================================

% Speaker of the following speech blocks (guest lecturer, change mid-lecture); \setSpeaker{} clears it.
% Shown faintly after the timestamp; without \setSpeaker nothing is shown.
\providecommand{\lecturespeaker}{}
\providecommand{\setSpeaker}[1]{\renewcommand{\lecturespeaker}{#1}}
\providecommand{\speakertag}{\ifdefempty{\lecturespeaker}{}{\quad{\mdseries\footnotesize\color{black!45}\faUser\ \ \lecturespeaker}}}

% 1. Clean Spoken Protocol
\newtcolorbox{speech}[1][]{
enhanced,
frame hidden,
colback=EnvSpokenBg,
dynamic shadow,
code={
  \ifinsideproof
    \tcbset{borderline west={3pt}{0pt}{EnvSpokenProof}}
  \else
    \tcbset{borderline west={3pt}{0pt}{EnvSpoken}}
  \fi
},
title={\sffamily\bfseries\textcolor{\ifinsideproof EnvSpokenProof\else EnvSpoken\fi}{\faMicrophone\ \ Mündliches Protokoll\ifblank{#1}{}{ [#1]}\speakertag}},
attach title to upper={\par\setlength{\parindent}{\standardparindent}\setlength{\parskip}{\customparskip}\vspace{1ex}},
breakable
}

% 2. AI Stroke-by-Stroke Note
\newtcolorbox{overview}[1][]{
enhanced,
colback=white,
boxrule=0.5pt,
dynamic shadow,
code={
  \ifinsideproof
    \tcbset{
      colframe=EnvNoteProof,
      boxed title style={colback=EnvNoteBg, colframe=EnvNoteProof, boxrule=0.5pt}
    }
  \else
    \tcbset{
      colframe=EnvNote,
      boxed title style={colback=EnvNoteBg, colframe=EnvNote, boxrule=0.5pt}
    }
  \fi
},
title={\sffamily\bfseries\faStickyNote\ \ #1},
coltitle=black,
attach boxed title to top left={xshift=10pt, yshift=-10pt},
top=14pt,
breakable
}

% 3. Best Teacher / Didactic Insight
\newtcolorbox{insight}[1][Didaktisches Konzept]{
enhanced,
arc=3mm,
frame hidden,
colback=EnvTeacherBg,
dynamic shadow,
code={
  \ifinsideproof
    \tcbset{borderline west={2pt}{0pt}{EnvTeacherProof}}
  \else
    \tcbset{borderline west={2pt}{0pt}{EnvTeacherBorder}}
  \fi
},
title={\sffamily\bfseries\textcolor{\ifinsideproof EnvTeacherProof\else EnvTeacherBorder\fi}{\textit{\faLightbulb\ \ Einblick:} #1}},
attach title to upper=\par\vspace{0.5ex},
breakable
}

% 4a. Modular Visual Blackboard Replication Box
\NewTColorBox{highlight}{ m o }{
enhanced,
colback=white,
colupper=black,
fonttitle=\bfseries\sffamily,
coltitle=black,
boxrule=1.5pt,
arc=4pt,
dynamic color shadow=#1!30,
code={
  \ifinsideproof
    \tcbset{
      colframe=#1!50!ProofTint,
      boxed title style={colback=white, colframe=#1!50!ProofTint, boxrule=0.5pt, top=2pt, bottom=2pt, left=4pt, right=4pt}
    }
  \else
    \tcbset{
      colframe=#1,
      boxed title style={colback=white, colframe=#1, boxrule=0.5pt, top=2pt, bottom=2pt, left=4pt, right=4pt}
    }
  \fi
},
IfValueTF={#2}{
  title={\faBox\ \ #2},
  attach boxed title to top left={xshift=10pt, yshift=-10pt},
  % the attached title reaches 10pt into the box: keep the first line clear of it (as overview)
  top=14pt
}{
  notitle
}
}

% 4b. Legacy Visual Blackboard Replication Boxes
\newcommand{\NewChalkBox}[3]{%
\newtcolorbox{#1}[1][Hervorgehobene Tafelformel]{%
enhanced,
colback=white,
fonttitle=\bfseries\sffamily,
coltitle=black,
title={\faBox\ \ ##1},
boxrule=1.5pt,
arc=4pt,
dynamic color shadow=#3,
attach boxed title to top left={xshift=10pt, yshift=-10pt},
top=14pt, % the attached title reaches 10pt into the box: keep the first line clear of it
code={
  \ifinsideproof
    \tcbset{
      colframe=#2!50!ProofTint,
      boxed title style={colback=white, colframe=#2!50!ProofTint, boxrule=0.5pt, top=2pt, bottom=2pt, left=4pt, right=4pt}
    }
  \else
    \tcbset{
      colframe=#2,
      boxed title style={colback=white, colframe=#2, boxrule=0.5pt, top=2pt, bottom=2pt, left=4pt, right=4pt}
    }
  \fi
}
}%
}

\NewChalkBox{orange-box}{BurntOrange}{BurntOrange!30}
\NewChalkBox{blue-box}{MidnightBlue}{MidnightBlue!30}
\NewChalkBox{dark-green-box}{ForestGreen}{ForestGreen!30}
\NewChalkBox{apple-green-box}{LimeGreen}{LimeGreen!30}
\NewChalkBox{yellow-box}{Goldenrod!90!black}{Goldenrod!30}
\NewChalkBox{violet-box}{Violet}{Violet!30}
\NewChalkBox{red-box}{BrickRed}{BrickRed!30}
\NewChalkBox{green-box}{Green}{Green!30}
\NewChalkBox{purple-box}{Purple}{Purple!30}
\NewChalkBox{gray-box}{Gray}{Gray!30}
\NewChalkBox{apricot-box}{Apricot}{Apricot!30}
\NewChalkBox{bittersweet-box}{Bittersweet}{Bittersweet!30}
\NewChalkBox{teal-blue-box}{TealBlue}{TealBlue!30}
\NewChalkBox{cyan-box}{Cyan}{Cyan!30}
\NewChalkBox{magenta-box}{Magenta}{Magenta!30}
\NewChalkBox{brown-box}{Brown}{Brown!30}
\NewChalkBox{pink-box}{Pink}{Pink!30}

\newcommand{\NewFormulaBox}[3]{%
\newtcolorbox{#1}[1][Hervorgehobene Tafelformel]{%
enhanced,
colback=white,
fonttitle=\bfseries\sffamily,
coltitle=black,
title={\faIcon{square-root-alt}\ \ ##1},
boxrule=1.5pt,
arc=4pt,
dynamic color shadow=#3,
attach boxed title to top left={xshift=10pt, yshift=-10pt},
top=14pt, % the attached title reaches 10pt into the box: keep the first line clear of it
code={
  \ifinsideproof
    \tcbset{
      colframe=#2!50!ProofTint,
      boxed title style={colback=white, colframe=#2!50!ProofTint, boxrule=0.5pt, top=2pt, bottom=2pt, left=4pt, right=4pt}
    }
  \else
    \tcbset{
      colframe=#2,
      boxed title style={colback=white, colframe=#2, boxrule=0.5pt, top=2pt, bottom=2pt, left=4pt, right=4pt}
    }
  \fi
}
}%
}

\NewFormulaBox{orange-formula}{BurntOrange}{BurntOrange!30}
\NewFormulaBox{blue-formula}{MidnightBlue}{MidnightBlue!30}
\NewFormulaBox{dark-green-formula}{ForestGreen}{ForestGreen!30}
\NewFormulaBox{apple-green-formula}{LimeGreen}{LimeGreen!30}
\NewFormulaBox{yellow-formula}{Goldenrod!90!black}{Goldenrod!30}
\NewFormulaBox{violet-formula}{Violet}{Violet!30}
\NewFormulaBox{red-formula}{BrickRed}{BrickRed!30}

% 5. Mathematical Setup (Stroke-by-Stroke)
\NewTColorBox{content}{ O{} }{
enhanced,
frame hidden,
colback=white,
opacityback=0, 
breakable,
before upper={\setlength{\parskip}{\customparskip}\setlength{\parindent}{0pt}},
code={
  \ifinsideproof
    \tcbset{borderline west={2pt}{0pt}{EnvMathProof}}
  \else
    \tcbset{borderline west={2pt}{0pt}{EnvMath!50}}
  \fi
  \ifblank{#1}{
    \tcbset{notitle}
  }{
    \tcbset{title={\sffamily\bfseries\textcolor{\ifinsideproof EnvMathProof\else EnvMath\fi}{\faChalkboard\ \ #1}}, 
    attach title to upper={\par\setlength{\parindent}{0pt}\setlength{\parskip}{\customparskip}\vspace{1ex}}}
  }
}
}

% 6. Meta-Note
\newtcolorbox{note}[1][Administrative Metanote]{
enhanced,
frame hidden,
colback=white,
opacityback=0,
code={
  \ifinsideproof
    \tcbset{borderline={0.5pt}{0pt}{EnvMetaProof, dashed}}
  \else
    \tcbset{borderline={0.5pt}{0pt}{EnvMeta, dashed}}
  \fi
},
title={\sffamily\bfseries\textcolor{\ifinsideproof EnvMetaProof\else EnvMeta\fi}{\faEye\ \ #1}},
attach title to upper=\par\vspace{1ex},
breakable
}

% 7. AI Observational Note
\newtcolorbox{aiNote}[1][KI-Notiz]{
enhanced,
colback=EnvAIBg,
dynamic shadow,
code={
  \ifinsideproof
    \tcbset{colframe=EnvAIProof}
  \else
    \tcbset{colframe=EnvAIBorder}
  \fi
},
fonttitle=\bfseries\sffamily,
coltitle=white,
title={\faRobot\ \ #1},
boxrule=1pt,
arc=3pt,
breakable
}

% 8. Student Question, Student Interaction, or Common Confusion.
\newtcolorbox{student}[1][Studentische Interaktion]{
enhanced,
frame hidden,
colback=EnvStudentBg,
dynamic shadow,
code={
  \ifinsideproof
    \tcbset{borderline west={3pt}{0pt}{EnvStudentProof}}
  \else
    \tcbset{borderline west={3pt}{0pt}{EnvStudent}}
  \fi
},
title={\sffamily\bfseries\textcolor{\ifinsideproof EnvStudentProof\else EnvStudent\fi}{\faUserGraduate\ \ #1}},
attach title to upper=\par\vspace{1ex},
breakable
}

% 9. Explanation of Steps
\newenvironment{steps}[1][Erklärung]{
\par\vspace{1ex}\noindent\textbf{\textcolor{ExplanationIconColor}{\faListUl\ \ #1:}}
}{
\par\vspace{1ex}
}

% 10. Subsection-like heading for proofs
\newenvironment{proof-subsection}[1]{
\par\vspace{1ex}\noindent
\tikz[baseline=(A.base)]\node[fill=ProofSubBg, text=ProofSubText, rounded corners=2pt, inner xsep=4pt, inner ysep=2pt, font=\sffamily\bfseries, drop shadow] (A) {#1};
\vspace{0.5ex}\par
}{
\par\vspace{1ex}
}

% 11. Dynamic Proof Wrappers
\makeatletter

% SHORT PROOF 
\renewenvironment{proof}[1][Beweis]{
    \par\addvspace{2ex plus 0.5ex minus 0.2ex}
    \insideprooftrue 
    \noindent\textcolor{ProofTint}{\faPenNib}\hspace{0.5em}%
    {\normalfont\sffamily\bfseries\color{ProofTint}#1}%
    \par\nobreak\vspace{1ex plus .2ex}\vspace{-\parskip}
    \@afterindentfalse\@afterheading
}{
    \par\vspace{-2ex}\hfill 
    \tikz[baseline=(Q.base)]\node[draw=ProofTint, fill=ProofQEDBg, line width=1pt, rounded corners=2pt, inner xsep=4pt, inner ysep=2pt, font=\sffamily\bfseries\small, text=ProofQEDText, drop shadow={shadow xshift=0.5ex, shadow yshift=-0.5ex, fill=ProofTint, opacity=0.4}] (Q) {Q.E.D.};
    \par\vspace{2ex}
}

% LONG PROOF
% With its own title, e.g. [Beweis von \cref{thm:x}], the proof gets a TOC entry and a PDF bookmark,
% one level below the heading it sits under; headings inside the proof are listed below it. With
% the default title (no option, or [Beweis]) it gets neither. The short proof never does.
\newcommand*{\proofWrapperDefaultTitle}{Beweis}
\let\proofwrapper@tocclose\relax
% Depth of the heading the proof sits under (0 chapter, 1 section, 2 subsection, 3 subsubsection),
% read off the counters: each heading resets the counter of the level below it.
\newcommand{\proofwrapper@depth}{%
  \ifnum\value{section}=\z@ 0%
  \else\ifnum\value{subsection}=\z@ 1%
  \else\ifnum\value{subsubsection}=\z@ 2%
  \else 3\fi\fi\fi}
\newcommand{\proofwrapper@tocentry}[1]{%
  \edef\proofwrapper@level{\proofwrapper@depth}%
  \phantomsection
  \addtocontents{toc}{\protect\proofwrappertocbegin{\proofwrapper@level}}%
  \edef\toclevel@proofwrapper{\the\numexpr\proofwrapper@level+1\relax}%
  \addcontentsline{toc}{proofwrapper}{\protect\proofwrappertoctitle{#1}}%
  % Bookmarks of the headings inside the proof become children of the proof's bookmark, by the
  % same rule as in the TOC: the first heading fixes the shift. The bookmark depth limit moves
  % with them, so that exactly the levels shown outside a proof are shown inside it.
  \gdef\proofwrapper@bmshift{-1}%
  \let\proofwrapper@addcontentsline\addcontentsline
  \def\addcontentsline##1##2##3{\proofwrapper@bmlevel{##2}\proofwrapper@addcontentsline{##1}{##2}{##3}}%
  \def\proofwrapper@tocclose{\addtocontents{toc}{\protect\proofwrappertocend}}%
}
\def\proofwrapper@baselevel@section{1}
\def\proofwrapper@baselevel@subsection{2}
\def\proofwrapper@baselevel@subsubsection{3}
\def\proofwrapper@baselevel@paragraph{4}
\def\proofwrapper@baselevel@subparagraph{5}
\newcommand{\proofwrapper@bmlevel}[1]{%
  \ifcsname proofwrapper@baselevel@#1\endcsname
    \ifnum\proofwrapper@bmshift<\z@
      \count@=\numexpr\proofwrapper@level+2-\csname proofwrapper@baselevel@#1\endcsname\relax
      \ifnum\count@<\z@ \count@=\z@ \fi
      \xdef\proofwrapper@bmshift{\the\count@}%
    \fi
    \expandafter\edef\csname toclevel@#1\endcsname
      {\the\numexpr\csname proofwrapper@baselevel@#1\endcsname+\proofwrapper@bmshift\relax}%
    \edef\Hy@bookmarksdepth{\the\numexpr\c@tocdepth+\proofwrapper@bmshift\relax}%
  \fi}
% PDF bookmarks cannot hold a \cref; spell it out from the label data cleveref writes to the aux
% file ([theorem][4][3]3.4 -> "Satz 3.4"), or "??" before the first run has written it.
\newcommand{\proofwrapper@pdfref}[1]{%
  \ifcsname r@#1@cref\endcsname
    \expandafter\@firstoftwo
  \else
    \expandafter\@secondoftwo
  \fi
  {\expandafter\expandafter\expandafter\proofwrapper@pdfref@i\csname r@#1@cref\endcsname\@nil}%
  {??}}
\def\proofwrapper@pdfref@i#1#2\@nil{\proofwrapper@pdfref@ii#1\@nil}
\def\proofwrapper@pdfref@ii[#1][#2][#3]#4\@nil{\csname Cref@#1@name\endcsname\space#4}
\pdfstringdefDisableCommands{%
  \let\proofwrappertoctitle\@firstofone
  \let\cref\proofwrapper@pdfref
  \let\Cref\proofwrapper@pdfref
}

\newenvironment{proofWrapper}[1][Beweis]{
    \par\addvspace{3.5ex plus 1ex minus .2ex}
    \insideprooftrue 
    \noindent
    \def\proofwrapper@title{#1}%
    \ifx\proofwrapper@title\proofWrapperDefaultTitle\else\proofwrapper@tocentry{#1}\fi
    \makebox[0pt][r]{\textcolor{ProofTint}{\faPenNib}\hspace{10pt}}%
    {\normalfont\sffamily\Large\bfseries\color{ProofTint}#1}%
    \par\nobreak\vspace{1.5ex plus .2ex}\vspace{-\parskip}
    \@afterindentfalse\@afterheading
}{
    \par\vspace{-4ex}\hfill 
    \tikz[baseline=(Q.base)]{
        \node[draw=ProofTint, fill=ProofQEDBg, line width=1pt, rounded corners=2pt, inner xsep=4pt, inner ysep=2pt, font=\sffamily\bfseries\small, text=ProofQEDText, drop shadow={shadow xshift=0.5ex, shadow yshift=-0.5ex, fill=ProofTint, opacity=0.4}] (Q) {Q.E.D.};
    }
    \par\vspace{2ex}
    \proofwrapper@tocclose
}
\makeatother

% 13. Lecture Break
\newtcolorbox{pause}[1][Vorlesungspause]{
enhanced,
frame hidden,
colback=white,
opacityback=0,
code={
  \ifinsideproof
    \tcbset{borderline north={0.5pt}{0pt}{EnvBreakProof, dashed}, borderline south={0.5pt}{0pt}{EnvBreakProof, dashed}}
  \else
    \tcbset{borderline north={0.5pt}{0pt}{EnvBreak, dashed}, borderline south={0.5pt}{0pt}{EnvBreak, dashed}}
  \fi
},
title={\sffamily\bfseries\centering\textcolor{\ifinsideproof EnvBreakProof\else EnvBreak\fi}{\faCoffee\ \ #1}},
attach title to upper=\par\vspace{1ex},
fontupper=\centering\itshape\color{\ifinsideproof EnvBreakProof\else EnvBreak\fi},
breakable
}

% ==========================================
% INVISIBLE LAYER & FALLBACK ENVIRONMENTS
% ==========================================
\NewEnviron{hidden}{}  % all AI scratchpads: plans, checkpoints, sketches

% Data Integrity Fallbacks
\newcommand{\NewFallbackBox}[1]{
\newtcolorbox{#1}[1][Fallback]{
enhanced,
colback=EnvAIBg,
dynamic shadow,
code={
  \ifinsideproof
    \tcbset{colframe=EnvAIFallbackProof}
  \else
    \tcbset{colframe=EnvAI}
  \fi
},
fonttitle=\sffamily\bfseries,
coltitle=white,
title={\faExclamationTriangle\ \ ##1},
breakable
}
}
\NewFallbackBox{ai-logical-gap}
\NewFallbackBox{ai-generation-loop-fallback}
\NewFallbackBox{ai-raw-ocr-fallback}
\NewFallbackBox{ai-phonetic-dump}
\NewFallbackBox{ai-modality-conflict}
\NewFallbackBox{ai-off-camera-state}
\NewFallbackBox{ai-async-board-update}
\NewFallbackBox{ai-kinetic-emphasis}

% Custom Lecture Chapter Command with Optional Semester/Week Number
\newif\ifrestorelecturesubcounters
\NewDocumentCommand{\lecturechapter}{ o m m m m }{%
  \restorelecturesubcountersfalse
  \IfValueT{#1}{%
    \ifblank{#1}{}{%
      \ifnum\numexpr#1\relax=\value{chapter}%
        \restorelecturesubcounterstrue
        \edef\TempSec{\the\value{section}}%
        \edef\TempThm{\the\value{theorem}}%
        \edef\TempEq{\the\value{equation}}%
        \edef\TempFig{\the\value{figure}}%
        \edef\TempTab{\the\value{table}}%
      \fi
      \setcounter{chapter}{\numexpr#1-1\relax}%
    }%
  }%
  \chapter[\texorpdfstring{--- \color{ThemePrimary!95!black}#5}{--- #5}]{{\Large\color{ChapSubtitleText}#2 #4} \texorpdfstring{\\[1ex]}{--- } {\LARGE #5}}%
  \ifrestorelecturesubcounters
    \setcounter{section}{\TempSec}%
    \setcounter{theorem}{\TempThm}%
    \setcounter{equation}{\TempEq}%
    \setcounter{figure}{\TempFig}%
    \setcounter{table}{\TempTab}%
  \fi
  \markboth{\chaptername\ \thechapter\ \ ---\ \ #2 #4\ \ ---\ \ #5}{}%
}

% Custom YouTube / Standalone Video Chapter Command:
% \youtubechapter[Date]{Creator / Channel}{Video Title}
\NewDocumentCommand{\youtubechapter}{ o m m }{%
  \IfValueTF{#1}{%
    \ifblank{#1}{%
      \chapter[\texorpdfstring{--- \color{ThemePrimary!95!black}#2: #3}{--- #2: #3}]{{\Large\color{ChapSubtitleText}#2} \texorpdfstring{\\[1ex]}{--- } {\LARGE #3}}%
      \markboth{\chaptername\ \thechapter\ \ ---\ \ #2\ \ ---\ \ #3}{}%
    }{%
      \chapter[\texorpdfstring{--- \color{ThemePrimary!95!black}#2 (#1): #3}{--- #2 (#1): #3}]{{\Large\color{ChapSubtitleText}#2\ ---\ #1} \texorpdfstring{\\[1ex]}{--- } {\LARGE #3}}%
      \markboth{\chaptername\ \thechapter\ \ ---\ \ #2\ (#1)\ \ ---\ \ #3}{}%
    }%
  }{%
    \chapter[\texorpdfstring{--- \color{ThemePrimary!95!black}#2: #3}{--- #2: #3}]{{\Large\color{ChapSubtitleText}#2} \texorpdfstring{\\[1ex]}{--- } {\LARGE #3}}%
    \markboth{\chaptername\ \thechapter\ \ ---\ \ #2\ \ ---\ \ #3}{}%
  }%
}

\titlespacing*{\chapter}{0pt}{-20pt}{40pt}

\titleformat{\chapter}[display]
{\normalfont\sffamily\Large\bfseries\color{ChapTitleText}}
{\raggedleft\large\color{ChapNumberText}\MakeUppercase{\chaptertitlename}\ \thechapter}
{1ex}
{\color{ChapLine}\rule{\textwidth}{1.5pt}\endgraf\nopagebreak\vspace{2ex}\raggedleft} 
[\endgraf\nopagebreak\vspace{1ex}\color{ChapLine}\rule{\textwidth}{1.5pt}]

% Extended Section Bars 
\newcommand{\sectionbar}{%
\tikz[baseline=0pt]{
\fill[\ifinsideproof SecBar1Proof\else SecBar1\fi] (-2.5pt, 0) rectangle (1.0pt, 0.72em);
\fill[\ifinsideproof SecBar2Proof\else SecBar2\fi] (2.0pt, 0) rectangle (3.5pt, 0.72em);
\fill[\ifinsideproof SecBar3Proof\else SecBar3\fi] (4.8pt, 0) rectangle (5.3pt, 0.72em);
\fill[\ifinsideproof SecBar4Proof\else SecBar4\fi] (6.6pt, 0) rectangle (8.1pt, 0.72em);
}%
\hspace{10pt}%
}

\newcommand{\subsectionbar}{%
\tikz[baseline=0pt]{
\fill[\ifinsideproof SubSecBar1Proof\else SubSecBar1\fi] (-2.5pt, 0) rectangle (0.3pt, 0.72em);
\shade[left color=\ifinsideproof SubSecBar1Proof\else SubSecBar1\fi, right color=\ifinsideproof SubSecBar2Proof\else SubSecBar2\fi] (1.3pt, 0) rectangle (2.0pt, 0.72em);
\fill[\ifinsideproof SubSecBar2Proof\else SubSecBar2\fi] (3.3pt, 0) rectangle (4.8pt, 0.72em);
}%
\hspace{10pt}%
}

\newcommand{\subsubsectionbar}{%
\tikz[baseline=0pt]{
\fill[\ifinsideproof SubSubSecBar1Proof\else SubSubSecBar1\fi] (-2.5pt, 0) rectangle (-1.0pt, 0.72em);
\fill[\ifinsideproof SubSubSecBar2Proof\else SubSubSecBar2\fi] (0,0) rectangle (2.2pt, 0.72em);
}%
\hspace{10pt}%
}

\titleformat{\section}{\normalfont\sffamily\Large\bfseries\color{\ifinsideproof SecTitleColorProof\else SecTitleColor\fi}}{}{0pt}{\phantomsection\makebox[0pt][r]{\sectionbar}}
\titleformat{\subsection}{\normalfont\sffamily\large\bfseries\color{\ifinsideproof SubSecTitleColorProof\else SubSecTitleColor\fi}}{}{0pt}{\phantomsection\makebox[0pt][r]{\subsectionbar}}
\titleformat{\subsubsection}{\normalfont\sffamily\normalsize\bfseries\color{\ifinsideproof SubSubSecTitleColorProof\else SubSubSecTitleColor\fi}}{}{0pt}{\phantomsection\makebox[0pt][r]{\subsubsectionbar}}
\titleformat{\paragraph}[runin]{\normalfont\normalsize\bfseries}{\theparagraph}{1em}{\phantomsection} % the class default, plus a link target on the heading (unnumbered, it has none)

\titlespacing*{\section}{0pt}{3.5ex plus 1ex minus .2ex}{1.5ex plus .2ex}
\titlespacing*{\subsection}{0pt}{3ex plus 1ex minus .2ex}{1.2ex plus .2ex}
\titlespacing*{\subsubsection}{0pt}{2.5ex plus 1ex minus .2ex}{1ex plus .2ex}


% =====================================================================
%  BACKWARD COMPATIBILITY (ACTIVE)
%
%  Lets .tex files written against the OLD vocabulary compile against this
%  preamble, so existing lecture output does not have to be migrated.
%  Delete this block once every .tex file uses the new names.
%
%  These are \let aliases rather than re-declared environments on purpose:
%  the alias inherits the target's argument signature and default title
%  exactly, so nothing here can drift out of sync with the definitions
%  above. Re-stating "[1][Erklaerung]" by hand would silently change the
%  rendered default the moment the real definition changed.
% =====================================================================
\makeatletter
\newcommand{\AliasEnv}[2]{%
  \expandafter\let\csname #1\expandafter\endcsname\csname #2\endcsname
  \expandafter\let\csname end#1\expandafter\endcsname\csname end#2\endcsname
}
\makeatother

\AliasEnv{spoken-clean}{speech}
\AliasEnv{math-stroke}{content}
\AliasEnv{student-interaction}{student}
\AliasEnv{explanation-of-steps}{steps}
\AliasEnv{didactic-insight}{insight}
\AliasEnv{meta-note}{note}
\AliasEnv{nice-box}{overview}
\AliasEnv{color-box}{highlight}
\AliasEnv{lecture-break}{pause}
\AliasEnv{ai-note}{aiNote}

%  CAVEAT -- the one thing aliasing cannot paper over.
%  The swap gave `proof` to the SHORT form (it is the common case, 63 uses
%  against 19). Old files use `proof` for the LONG form and `short-proof`
%  for the short one, so `proof` means opposite things in the two
%  vocabularies and no single preamble can serve both.
%  `short-proof` is aliased to the new `proof`, which is correct. Old
%  `\begin{proof}` therefore renders with the short-proof heading --
%  slightly smaller, no hanging icon. It still compiles, still reads
%  "Beweis", still emits Q.E.D. Cosmetic only, and it disappears as soon
%  as the file is migrated.
\AliasEnv{short-proof}{proof}

%  Scratchpads: all of these swallow their body, so an empty definition is
%  exactly equivalent to aliasing them to `hidden`.
\NewEnviron{ai-tikz-planner-invisible-content}{}
\NewEnviron{ai-proof-skeleton-invisible-content}{}
\NewEnviron{ai-global-state-checkpoint-invisible-content}{}
\NewEnviron{ai-retroactive-patch}[1][]{}
\NewEnviron{ai-global-state-checkpoint}{}

%  NOT merged into `hidden`, deliberately: the fallback boxes below are
%  \newtcolorbox, not \NewEnviron -- they RENDER, as red warning panels. They
%  exist to surface OCR failures, off-camera state and modality conflicts for a
%  human to see. Making them invisible would silently swallow exactly the
%  content they were created to raise. They keep their own names and stay
%  visible: ai-logical-gap, ai-generation-loop-fallback, ai-raw-ocr-fallback,
%  ai-phonetic-dump, ai-modality-conflict, ai-off-camera-state,
%  ai-async-board-update, ai-kinetic-emphasis.
